\begin{abstract}
We measure how nine Korean sentiment pipelines change from 200 labeled reviews to each dataset's full cleaned training pool. The main benchmark covers NSMC movie reviews, Naver Shopping reviews, and Korean Steam reviews, with four TF-IDF classifiers and five neural models learned from scratch. Every final condition uses three fixed seeds and shared nested subsets. The best neural mean exceeded the best TF-IDF mean by at least one macro-F1 point, and kept that margin, at \Result{part1_crossover_nsmc} labels on NSMC. The criterion was \Result{part1_crossover_shopping} on Shopping and \Result{part1_crossover_steam} on Steam. At full data, the best neural gains over the best TF-IDF pipeline were \Result{story_full_gap_nsmc}, \Result{story_full_gap_shopping}, and \Result{story_full_gap_steam} percentage points, respectively, while mean fit time rose by factors of \Result{story_fit_ratio_nsmc}, \Result{story_fit_ratio_shopping}, and \Result{story_fit_ratio_steam}. A bonus extension adds two labeled source domains to a low-label target. Pooling improved the target-only Transformer at 200 labels in all rotations, but standard SupCon and DB-CDSupCon produced mixed incremental effects. Target-only crossovers were not observed within 200--2,000 labels, so new target-label savings could not be estimated. The practical result is scale-dependent: linear TF-IDF is a strong low-cost default, while learned-from-scratch neural models become competitive only near the largest labeled budgets.
\end{abstract}
